
\documentclass{article}
\usepackage{colortbl}
\usepackage{makecell}
\usepackage{multirow}
\usepackage{supertabular}

\begin{document}

\newcounter{utterance}

\centering \large Interaction Transcript for game `codenames', experiment `opponent\_difficult', episode 2 with Qwen3.6{-}35B{-}A3B{-}FP8{-}without{-}reasoning.
\vspace{24pt}

{ \footnotesize  \setcounter{utterance}{1}
\setlength{\tabcolsep}{0pt}
\begin{supertabular}{c@{$\;$}|p{.15\linewidth}@{}p{.15\linewidth}p{.15\linewidth}p{.15\linewidth}p{.15\linewidth}p{.15\linewidth}}
    \# & $\;$A & \multicolumn{4}{c}{Game Master} & $\;\:$B\\
    \hline

    \theutterance \stepcounter{utterance}  
    & & & \multicolumn{2}{p{0.3\linewidth}}{
        \cellcolor[rgb]{0.9,0.9,0.9}{
            \makecell[{{p{\linewidth}}}]{
                \texttt{\tiny{[GM$|$GM]}}
                {'hidden': {'team': ['konj', 'veleposlanstvo', 'centar', 'metvica', 'vrijeme', 'bič', 'točka', 'oko', 'crkva'], 'innocent': ['prekidač', 'pita', 'vlak', 'glavni grad', 'bermuda', 'trgovina', 'agent'], 'opponent': ['palac', 'duša', 'otrov', 'francuska', 'krtica', 'opera', 'litica', 'ekran'], 'assassin': ['crv']}, 'revealed': {'team': {'team': [], 'innocent': [], 'opponent': [], 'assassin': []}, 'opponent': {'team': [], 'innocent': [], 'opponent': [], 'assassin': []}}}
            }
        }
    }
    & & \\ \\

    \theutterance \stepcounter{utterance}  
    & & \multicolumn{4}{p{0.6\linewidth}}{
        \cellcolor[rgb]{0.9,0.9,0.9}{
            \makecell[{{p{\linewidth}}}]{
                \texttt{\tiny{[P1$\langle$GM]}}
                \texttt{Navedite jednu jednorječnu natuknicu koja se odnosi na neke riječi vašeg tima s popisa u nastavku. Možete odabrati na koliko i na koje se riječi vaša natuknica odnosi.} \\
\texttt{Cilj je da vaš partner pogodi sve riječi vašeg tima na temelju natuknica koje mu dajete, prije nego što drugi tim otkrije sve svoje. Ako vaš partner pogodi riječ 'protivnika' ili riječ 'mamac', njegov potez pogađanja završava i suprotni tim dobiva priliku odigrati potez. Ako vaš partner pogodi riječ 'ubojica', igra završava i vaš tim odmah gubi.} \\
\texttt{Što se vaša natuknica odnosi na više riječi, to brže možete pobijediti, ali asocijacije mogu biti teže za pogoditi vašem partneru, zato pažljivo odaberite svoju natuknicu. Natuknica mora biti semantički povezana s ciljanim riječima; ne smije biti jedna od riječi na popisima niti sadržavati njihove dijelove.} \\
\\ 
\texttt{Uvijek navedite svoju jednorječnu natuknicu i svoj popis povezanih ciljanih riječi odvojen zarezima u sljedećem formatu i neka vaši odgovori budu što kraći; nikada ne uključujte nikakav drugi tekst osim onoga što je potrebno u ovom obrascu:} \\
\\ 
\texttt{NATUKNICA: <RIJEČ>} \\
\texttt{CILJEVI: <CILJEVI>} \\
\\ 
\texttt{Riječi vašeg tima su: konj, veleposlanstvo, centar, metvica, vrijeme, bič, točka, oko, crkva.} \\
\texttt{Riječi protivničkog tima su: palac, duša, otrov, francuska, krtica, opera, litica, ekran.} \\
\texttt{Riječi mamci su: prekidač, pita, vlak, glavni grad, bermuda, trgovina, agent.} \\
\texttt{Riječi ubojice su: crv.} \\
            }
        }
    }
    & \\ \\

    \theutterance \stepcounter{utterance}  
    & \multicolumn{4}{p{0.6\linewidth}}{
        \cellcolor[rgb]{0.9,0.9,0.9}{
            \makecell[{{p{\linewidth}}}]{
                \texttt{\tiny{[P1$\rangle$GM]}}
                \texttt{NATUKNICA: životinja} \\
\texttt{CILJEVI: konj, krtica} \\
            }
        }
    }
    & & \\ \\

    \theutterance \stepcounter{utterance}  
    & & & \multicolumn{2}{p{0.3\linewidth}}{
        \cellcolor[rgb]{0.9,0.9,0.9}{
            \makecell[{{p{\linewidth}}}]{
                \texttt{\tiny{[GM$|$GM]}}
                \texttt{životinja} \\
            }
        }
    }
    & & \\ \\

    \theutterance \stepcounter{utterance}  
    & & & \multicolumn{2}{p{0.3\linewidth}}{
        \cellcolor[rgb]{0.9,0.9,0.9}{
            \makecell[{{p{\linewidth}}}]{
                \texttt{\tiny{[GM$|$GM]}}
                ['konj', 'krtica']
            }
        }
    }
    & & \\ \\

    \theutterance \stepcounter{utterance}  
    & & & \multicolumn{2}{p{0.3\linewidth}}{
        \cellcolor[rgb]{0.9,0.9,0.9}{
            \makecell[{{p{\linewidth}}}]{
                \texttt{\tiny{[GM$|$GM]}}
                {'word': 'konj', 'assignment': 'team'}
            }
        }
    }
    & & \\ \\

    \theutterance \stepcounter{utterance}  
    & & & \multicolumn{2}{p{0.3\linewidth}}{
        \cellcolor[rgb]{0.9,0.9,0.9}{
            \makecell[{{p{\linewidth}}}]{
                \texttt{\tiny{[GM$|$GM]}}
                {'word': 'krtica', 'assignment': 'opponent'}
            }
        }
    }
    & & \\ \\

    \theutterance \stepcounter{utterance}  
    & & \multicolumn{4}{p{0.6\linewidth}}{
        \cellcolor[rgb]{0.9,0.9,0.9}{
            \makecell[{{p{\linewidth}}}]{
                \texttt{\tiny{[GM$\rangle$P2]}}
                \texttt{Navedite popis do 2 riječi odvojenih zarezom s donjeg popisa koje su najbolje povezane ili najuže povezane s riječju 'životinja'. Uvijek započnite svoj popis pogađanja s POGAĐANJE:  i ne uključujte nikakav drugi tekst u svoj odgovor.} \\
\\ 
\texttt{opera, francuska, veleposlanstvo, bermuda, točka, bič, krtica, oko, prekidač, palac, centar, trgovina, vlak, agent, vrijeme, ekran, pita, crv, glavni grad, konj, metvica, crkva, otrov, duša, litica} \\
            }
        }
    }
    & \\ \\

    \theutterance \stepcounter{utterance}  
    & & & \multicolumn{4}{p{0.6\linewidth}}{
        \cellcolor[rgb]{0.9,0.9,0.9}{
            \makecell[{{p{\linewidth}}}]{
                \texttt{\tiny{[GM$\langle$P2]}}
                \texttt{POGAĐANJE: krtica, konj} \\
            }
        }
    }
     \\ \\

    \theutterance \stepcounter{utterance}  
    & & & \multicolumn{2}{p{0.3\linewidth}}{
        \cellcolor[rgb]{0.9,0.9,0.9}{
            \makecell[{{p{\linewidth}}}]{
                \texttt{\tiny{[GM$|$GM]}}
                ['krtica', 'konj']
            }
        }
    }
    & & \\ \\

    \theutterance \stepcounter{utterance}  
    & & & \multicolumn{2}{p{0.3\linewidth}}{
        \cellcolor[rgb]{0.9,0.9,0.9}{
            \makecell[{{p{\linewidth}}}]{
                \texttt{\tiny{[GM$|$GM]}}
                {'word': 'krtica', 'assignment': 'opponent'}
            }
        }
    }
    & & \\ \\

    \theutterance \stepcounter{utterance}  
    & & & \multicolumn{2}{p{0.3\linewidth}}{
        \cellcolor[rgb]{0.9,0.9,0.9}{
            \makecell[{{p{\linewidth}}}]{
                \texttt{\tiny{[GM$|$GM]}}
                {'word': 'krtica', 'assignment': 'opponent'}
            }
        }
    }
    & & \\ \\

    \theutterance \stepcounter{utterance}  
    & & & \multicolumn{2}{p{0.3\linewidth}}{
        \cellcolor[rgb]{0.9,0.9,0.9}{
            \makecell[{{p{\linewidth}}}]{
                \texttt{\tiny{[GM$|$GM]}}
                \texttt{krtica} \\
            }
        }
    }
    & & \\ \\

    \theutterance \stepcounter{utterance}  
    & & & \multicolumn{2}{p{0.3\linewidth}}{
        \cellcolor[rgb]{0.9,0.9,0.9}{
            \makecell[{{p{\linewidth}}}]{
                \texttt{\tiny{[GM$|$GM]}}
                {'word': 'palac', 'assignment': 'opponent'}
            }
        }
    }
    & & \\ \\

    \theutterance \stepcounter{utterance}  
    & & & \multicolumn{2}{p{0.3\linewidth}}{
        \cellcolor[rgb]{0.9,0.9,0.9}{
            \makecell[{{p{\linewidth}}}]{
                \texttt{\tiny{[GM$|$GM]}}
                {'word': 'duša', 'assignment': 'opponent'}
            }
        }
    }
    & & \\ \\

    \theutterance \stepcounter{utterance}  
    & & & \multicolumn{2}{p{0.3\linewidth}}{
        \cellcolor[rgb]{0.9,0.9,0.9}{
            \makecell[{{p{\linewidth}}}]{
                \texttt{\tiny{[GM$|$GM]}}
                {'hidden': {'team': ['konj', 'veleposlanstvo', 'centar', 'metvica', 'vrijeme', 'bič', 'točka', 'oko', 'crkva'], 'innocent': ['prekidač', 'pita', 'vlak', 'glavni grad', 'bermuda', 'trgovina', 'agent'], 'opponent': ['otrov', 'francuska', 'opera', 'litica', 'ekran'], 'assassin': ['crv']}, 'revealed': {'team': {'team': [], 'innocent': [], 'opponent': ['krtica'], 'assassin': []}, 'opponent': {'team': [], 'innocent': [], 'opponent': ['palac', 'duša'], 'assassin': []}}}
            }
        }
    }
    & & \\ \\

    \theutterance \stepcounter{utterance}  
    & & \multicolumn{4}{p{0.6\linewidth}}{
        \cellcolor[rgb]{0.9,0.9,0.9}{
            \makecell[{{p{\linewidth}}}]{
                \texttt{\tiny{[P1$\langle$GM]}}
                \texttt{Riječ krtica je pogođena, ali je to protivnička riječ. Poteza tvoga suigrača završila se tamo.} \\
\texttt{Sada daj još jednu natuknicu koja se odnosi na neke od preostalih riječi tvog tima i popis povezanih ciljnih riječi. Ne zaboravi započeti svoju natuknicu s 'NATUKNICA: ', staviti novi redak i započeti svoj popis ciljnih riječi odvojenih zarezima s 'CILJEVI: '. Napomena: neke riječi su uklonjene s popisa u usporedbi s prethodnim zahtjevima.} \\
\\ 
\texttt{Tvoje preostale timske riječi su: konj, veleposlanstvo, centar, metvica, vrijeme, bič, točka, oko, crkva.} \\
\texttt{Preostale riječi za tvog protivnika su: palac, duša, otrov, francuska, opera, litica, ekran.} \\
\texttt{Preostale riječi ometača su: prekidač, pita, vlak, glavni grad, bermuda, trgovina, agent.} \\
\texttt{Preostale ubilačke riječi su: crv.} \\
            }
        }
    }
    & \\ \\

    \theutterance \stepcounter{utterance}  
    & \multicolumn{4}{p{0.6\linewidth}}{
        \cellcolor[rgb]{0.9,0.9,0.9}{
            \makecell[{{p{\linewidth}}}]{
                \texttt{\tiny{[P1$\rangle$GM]}}
                \texttt{NATUKNICA: zgrada} \\
\texttt{CILJEVI: veleposlanstvo, crkva} \\
            }
        }
    }
    & & \\ \\

    \theutterance \stepcounter{utterance}  
    & & & \multicolumn{2}{p{0.3\linewidth}}{
        \cellcolor[rgb]{0.9,0.9,0.9}{
            \makecell[{{p{\linewidth}}}]{
                \texttt{\tiny{[GM$|$GM]}}
                \texttt{zgrada} \\
            }
        }
    }
    & & \\ \\

    \theutterance \stepcounter{utterance}  
    & & & \multicolumn{2}{p{0.3\linewidth}}{
        \cellcolor[rgb]{0.9,0.9,0.9}{
            \makecell[{{p{\linewidth}}}]{
                \texttt{\tiny{[GM$|$GM]}}
                ['veleposlanstvo', 'crkva']
            }
        }
    }
    & & \\ \\

    \theutterance \stepcounter{utterance}  
    & & & \multicolumn{2}{p{0.3\linewidth}}{
        \cellcolor[rgb]{0.9,0.9,0.9}{
            \makecell[{{p{\linewidth}}}]{
                \texttt{\tiny{[GM$|$GM]}}
                {'word': 'veleposlanstvo', 'assignment': 'team'}
            }
        }
    }
    & & \\ \\

    \theutterance \stepcounter{utterance}  
    & & & \multicolumn{2}{p{0.3\linewidth}}{
        \cellcolor[rgb]{0.9,0.9,0.9}{
            \makecell[{{p{\linewidth}}}]{
                \texttt{\tiny{[GM$|$GM]}}
                {'word': 'crkva', 'assignment': 'team'}
            }
        }
    }
    & & \\ \\

    \theutterance \stepcounter{utterance}  
    & & \multicolumn{4}{p{0.6\linewidth}}{
        \cellcolor[rgb]{0.9,0.9,0.9}{
            \makecell[{{p{\linewidth}}}]{
                \texttt{\tiny{[GM$\rangle$P2]}}
                \texttt{Riječ krtica je pogođena, ali je to protivnička riječ. Tvoja poteza je završena tamo.} \\
\texttt{Sada navedite još jedan popis odvojen zarezima od najmanje 1 i najviše 2 riječi iz sljedećeg popisa riječi koje se najbolje odnose ili su najuže povezane s riječju 'zgrada'. Ne zaboravite započeti svoj odgovor s POGAĐANJE: . Napomena: neke su riječi uklonjene s popisa u odnosu na prethodne zahtjeve.} \\
\\ 
\texttt{opera, francuska, veleposlanstvo, bermuda, točka, bič, oko, prekidač, centar, trgovina, vlak, agent, vrijeme, ekran, pita, crv, glavni grad, konj, metvica, crkva, otrov, litica} \\
            }
        }
    }
    & \\ \\

    \theutterance \stepcounter{utterance}  
    & & & \multicolumn{4}{p{0.6\linewidth}}{
        \cellcolor[rgb]{0.9,0.9,0.9}{
            \makecell[{{p{\linewidth}}}]{
                \texttt{\tiny{[GM$\langle$P2]}}
                \texttt{POGAĐANJE: opera, crkva} \\
            }
        }
    }
     \\ \\

    \theutterance \stepcounter{utterance}  
    & & & \multicolumn{2}{p{0.3\linewidth}}{
        \cellcolor[rgb]{0.9,0.9,0.9}{
            \makecell[{{p{\linewidth}}}]{
                \texttt{\tiny{[GM$|$GM]}}
                ['opera', 'crkva']
            }
        }
    }
    & & \\ \\

    \theutterance \stepcounter{utterance}  
    & & & \multicolumn{2}{p{0.3\linewidth}}{
        \cellcolor[rgb]{0.9,0.9,0.9}{
            \makecell[{{p{\linewidth}}}]{
                \texttt{\tiny{[GM$|$GM]}}
                {'word': 'opera', 'assignment': 'opponent'}
            }
        }
    }
    & & \\ \\

    \theutterance \stepcounter{utterance}  
    & & & \multicolumn{2}{p{0.3\linewidth}}{
        \cellcolor[rgb]{0.9,0.9,0.9}{
            \makecell[{{p{\linewidth}}}]{
                \texttt{\tiny{[GM$|$GM]}}
                \texttt{opera} \\
            }
        }
    }
    & & \\ \\

    \theutterance \stepcounter{utterance}  
    & & & \multicolumn{2}{p{0.3\linewidth}}{
        \cellcolor[rgb]{0.9,0.9,0.9}{
            \makecell[{{p{\linewidth}}}]{
                \texttt{\tiny{[GM$|$GM]}}
                {'word': 'litica', 'assignment': 'opponent'}
            }
        }
    }
    & & \\ \\

    \theutterance \stepcounter{utterance}  
    & & & \multicolumn{2}{p{0.3\linewidth}}{
        \cellcolor[rgb]{0.9,0.9,0.9}{
            \makecell[{{p{\linewidth}}}]{
                \texttt{\tiny{[GM$|$GM]}}
                {'word': 'otrov', 'assignment': 'opponent'}
            }
        }
    }
    & & \\ \\

    \theutterance \stepcounter{utterance}  
    & & & \multicolumn{2}{p{0.3\linewidth}}{
        \cellcolor[rgb]{0.9,0.9,0.9}{
            \makecell[{{p{\linewidth}}}]{
                \texttt{\tiny{[GM$|$GM]}}
                {'hidden': {'team': ['konj', 'veleposlanstvo', 'centar', 'metvica', 'vrijeme', 'bič', 'točka', 'oko', 'crkva'], 'innocent': ['prekidač', 'pita', 'vlak', 'glavni grad', 'bermuda', 'trgovina', 'agent'], 'opponent': ['francuska', 'ekran'], 'assassin': ['crv']}, 'revealed': {'team': {'team': [], 'innocent': [], 'opponent': ['krtica', 'opera'], 'assassin': []}, 'opponent': {'team': [], 'innocent': [], 'opponent': ['palac', 'duša', 'litica', 'otrov'], 'assassin': []}}}
            }
        }
    }
    & & \\ \\

    \theutterance \stepcounter{utterance}  
    & & \multicolumn{4}{p{0.6\linewidth}}{
        \cellcolor[rgb]{0.9,0.9,0.9}{
            \makecell[{{p{\linewidth}}}]{
                \texttt{\tiny{[P1$\langle$GM]}}
                \texttt{Riječ opera je pogođena, ali je to protivnička riječ. Poteza tvoga suigrača završila se tamo.} \\
\texttt{Sada daj još jednu natuknicu koja se odnosi na neke od preostalih riječi tvog tima i popis povezanih ciljnih riječi. Ne zaboravi započeti svoju natuknicu s 'NATUKNICA: ', staviti novi redak i započeti svoj popis ciljnih riječi odvojenih zarezima s 'CILJEVI: '. Napomena: neke riječi su uklonjene s popisa u usporedbi s prethodnim zahtjevima.} \\
\\ 
\texttt{Tvoje preostale timske riječi su: konj, veleposlanstvo, centar, metvica, vrijeme, bič, točka, oko, crkva.} \\
\texttt{Preostale riječi za tvog protivnika su: otrov, francuska, litica, ekran.} \\
\texttt{Preostale riječi ometača su: prekidač, pita, vlak, glavni grad, bermuda, trgovina, agent.} \\
\texttt{Preostale ubilačke riječi su: crv.} \\
            }
        }
    }
    & \\ \\

    \theutterance \stepcounter{utterance}  
    & \multicolumn{4}{p{0.6\linewidth}}{
        \cellcolor[rgb]{0.9,0.9,0.9}{
            \makecell[{{p{\linewidth}}}]{
                \texttt{\tiny{[P1$\rangle$GM]}}
                \texttt{NATUKNICA: dio tijela} \\
\texttt{CILJEVI: oko, bič} \\
            }
        }
    }
    & & \\ \\

    \theutterance \stepcounter{utterance}  
    & & & \multicolumn{2}{p{0.3\linewidth}}{
        \cellcolor[rgb]{0.9,0.9,0.9}{
            \makecell[{{p{\linewidth}}}]{
                \texttt{\tiny{[GM$|$GM]}}
                {'player': 'Cluegiver', 'type': 'clue contains spaces', 'utterance': 'NATUKNICA: dio tijela\nCILJEVI: oko, bič', 'message': "Clue 'dio tijela' contains spaces and thus is not a single word.", 'clue': 'dio tijela'}
            }
        }
    }
    & & \\ \\

\end{supertabular}
}

\end{document}
